\chapter{Verification Methodology}\label{ch10:testing}

A kernel whose entire purpose is to make invalid inferences unwritable must
hold itself to a correspondingly explicit standard of evidence. This
chapter documents how the TKS v8.0 prototype is tested: the
bounded-exhaustive philosophy that governs the suite design, the four test
modules and the exact coverage they achieve, the boundary regression that
pins the artifact's claim discipline into the suite itself, the Gate~1
verifier that binds hashes to test runs, and --- given equal weight
throughout --- a precise statement of what a green suite does \emph{not}
establish. The numbers in this chapter are the numbers of the validation
snapshot of 2026-08-04 recorded in the prototype manifest: result
\textsc{pass}, 20 tests run, with bounded exhaustive coverage of
$19{,}683$ associativity triples, $289$ endomaps, $499$ typed maps, and
$1{,}678$ lift-functor composition pairs. No number below is a projection
or an extrapolation; each is either quoted from the manifest or derived
arithmetically and marked as a derivation.

\section{The bounded-exhaustive philosophy}\label{ch10:sec:philosophy}

\subsection{Small scope as a working hypothesis}

The suite's design principle is \emph{exhaustion within an explicit bound}:
rather than sampling large instances at random, enumerate \emph{every}
instance up to a small size and check the property on all of them. The
methodological warrant is the small-scope hypothesis articulated by
Jackson in the context of the Alloy analyzer: a high proportion of bugs in
declarative models and structured code can be exhibited by counterexamples
of small size, so exhaustively checking all small instances is a
disproportionately effective use of finite checking budget
\cite{jackson}. The hypothesis is empirical and defeasible --- it is a
claim about where bugs tend to live, not a theorem about where they must
--- and the suite treats it as exactly that: a rational allocation policy
for executable evidence, never a substitute for proof.

For this kernel the hypothesis has unusually strong purchase, for a
structural reason worth spelling out. The properties under test ---
associativity of index-chasing composition, neutrality of identities,
idempotence of factorial powers, functoriality of the lift --- are
\emph{uniform} in carrier size: the code that handles a three-element
carrier is the same code, on the same control path, that handles a
three-hundred-element carrier. There are no size-dependent branches, no
buffering thresholds, no algorithmic mode switches anywhere in the three
modules. A defect in the composition formula, the powering loop, or the
matrix builder would therefore almost certainly manifest at the smallest
size at which the relevant structure exists --- collisions need two
colliding elements, associativity needs three composable arrows,
nontrivial cycles need two points. Exhausting small scopes covers every
\emph{shape} of instance even though it cannot cover every \emph{size},
and for uniform code, shape is where the bugs are.

\subsection{Contrast: sampled property-based testing}

The natural alternative is property-based random testing in the QuickCheck
tradition \cite{quickcheck}: state the law as an executable predicate,
draw pseudorandom instances from a generator distribution, and shrink any
failure to a minimal counterexample. The two disciplines share the crucial
first step --- both force the specification into executable form --- and
differ in their quantifier. A QuickCheck-style run establishes ``the
property held on $n$ samples drawn from this distribution''; a
bounded-exhaustive run establishes ``the property holds on \emph{all}
instances within this bound.'' The former statement is probabilistic and
distribution-relative: its evidential weight depends on how well the
generator's measure covers the shapes where defects hide, and a generator
that rarely produces, say, empty carriers or maximal-collision mappings
can pass indefinitely against a defect it never exercises. The latter
statement is categorical within its bound and requires no distributional
argument at all --- but it is silent beyond the bound.

The kernel's suite chooses exhaustion for three reasons. First,
determinism: a random suite's failures depend on seeds, and this artifact's
entire identity discipline is a war on interpreter-state dependence; a
test suite whose coverage varied by seed would be tonally at odds with the
code it checks, and a replayable suite is a precondition for the
hash-bound verification protocol of Section~\ref{ch10:sec:gate1}. Second,
countability: the instance spaces at small scope are so small (hundreds to
tens of thousands, computed exactly below) that exhaustion is cheaper than
any statistically meaningful sampling campaign --- when the whole space
costs less than a second, sampling from it is strictly dominated. Third,
auditability: ``all $1{,}678$ composable pairs through size three'' is a
claim a reviewer can recount from first principles, as we do below,
whereas a sampling claim can only be re-audited by rerunning the sampler.
None of this is a criticism of the sampled discipline in its home
territory --- vast or infinite instance spaces where exhaustion is
impossible --- but that territory is precisely what the kernel's finite,
small-scope design avoids.

\begin{boundary}\label{ch10:bdy:scope}
Exhaustive-within-bounds is still bounded. Every coverage claim in this
chapter carries an explicit scope --- carriers through size three or four
--- and establishes nothing, even empirically, about behavior beyond that
scope. The bridge from bounded evidence to universal statements is
mathematical proof (Chapters~4--8) and, where present, kernel-checked
formalization (Chapter~11); the suite's role is to tie the running code to
the proved mathematics on the bounded family, not to extend the
mathematics.
\end{boundary}

\subsection{The oracle problem}

Every test needs an oracle: an independent source of the expected answer.
The suites resolve the oracle problem differently per family, and the
differences are methodologically instructive.

For the \emph{law families} --- associativity, identity neutrality,
functoriality --- the oracle is internal to the law itself: the test
computes both sides of an equation through the implementation and asserts
their equality. This is the cheapest kind of oracle and also the weakest
in one specific sense: a defect that breaks both sides \emph{identically}
--- say, a composition routine that consistently reversed its arguments
--- can satisfy an equational assertion while computing the wrong
function. The suites compensate by anchoring the equations to
independently constructed values at their edges: identities are built by
an independent constructor (\texttt{identity}), not by composing anything,
so the unit laws compare a computed composite against a value the
composition code never touched; and the lift functor law compares a
matrix \emph{product}, computed by explicit integer summation in the test,
against a matrix built directly from the composed arrow --- two
genuinely different computational routes to the same object.

For the \emph{theorem families} --- factorial idempotence, the norm
consequences --- the oracle is the informal proof: the test asserts the
proved right-hand side (idempotence of $f^{|X|!}$; trace equal to $|X|$;
rank equal to image size) against implementation output. Here the
mathematics is the independent authority, and the test would catch either
an implementation defect or --- more interestingly --- a misremembered
theorem, since the asserted value is written into the test from the
proof, not read from the code.

For the \emph{fixture families} --- the two-state noncontractivity
example, the hand-worked closure catalogs --- the oracle is hand
computation, of exactly the kind Appendix~A performs in print. These are
the strongest oracles in the suite (fully independent of both the
implementation and any single theorem) and necessarily the narrowest.

\subsection{Determinism of the suite itself}

The suite contains no randomness: no seeds, no time-dependent values, no
reliance on dictionary iteration order in any assertion. Enumeration of
the exhaustive families is by nested iteration over mapping tuples in
lexicographic order --- itself a replayable sequence --- so not only the
verdict but the complete execution trace of the suite is identical on
every conforming interpreter. This property is not merely aesthetic
consistency with the kernel's determinism discipline. It is what makes
the Gate~1 verifier's conjunction meaningful: ``these bytes pass these
tests'' is a reproducible fact only if ``these tests'' denotes one fixed
computation, and a suite with ambient nondeterminism would smuggle
interpreter state back into the artifact's evidence base through the
test channel after the implementation had so carefully evicted it.

\section{The four suites}\label{ch10:sec:suites}

The test tree \texttt{tests/tks\_v8\_core} contains four modules, each
hash-bound in the manifest with a declared role
(Chapter~9, Section on manifest binding): \texttt{test\_category.py}
(category and power regressions, 5{,}588 bytes),
\texttt{test\_deterministic\_lift.py} (lift structure and functor-law
regressions, 9{,}142 bytes), \texttt{test\_normal\_forms.py} (normal-form
correctness and resource regressions, 5{,}672 bytes), and
\texttt{test\_manifest\_boundaries.py} (claim and integration boundary
regression, 1{,}965 bytes). Together they comprise the 20 tests of the
snapshot. The suites run under Python's standard \texttt{unittest}
discovery --- \texttt{python -m unittest discover -s tests/tks\_v8\_core
-v} with \texttt{src} on \texttt{PYTHONPATH} --- with no third-party test
framework, no fixtures requiring network or filesystem state, and no
randomness anywhere: every run of the suite is the same run.

\subsection{Category and power regressions}\label{ch10:sec:cat-suite}

The category suite carries the two largest exhaustive families.

\emph{Associativity on a three-point carrier: $19{,}683$ triples.} Fix one
three-element carrier $X$. The endo-arrows on $X$ are the total functions
from a three-element set to itself, of which there are $3^3 = 27$; every
ordered triple $(f, g, h)$ of them is composable in both associations, and
there are $27^3 = (3^3)^3 = 19{,}683$ such triples. The suite composes
each triple both ways --- $(h \circ g) \circ f$ and $h \circ (g \circ f)$
--- and asserts extensional equality of the results. This family checks
the associativity component of the category laws (Chapter~4) on every
possible shape of collision, fixed point, and cycle structure that three
points admit, which includes every way two composed collisions can
interact. Identity neutrality is checked on typed families as well as
endo families, so the unit laws are not confined to the single-carrier
world.

\emph{Endomaps through size four: $289$ maps.} For each carrier size
$k \in \{0, 1, 2, 3, 4\}$ the endomaps on a $k$-element carrier number
$k^k$ (with the empty case counting the one empty map, $0^0 = 1$):

\[
\sum_{k=0}^{4} k^k \;=\; 1 + 1 + 4 + 27 + 256 \;=\; 289 ,
\]

a derivation consistent with the manifest's recorded count. For every one
of these $289$ endomaps the suite computes the factorial power
$f^{|X|!}$ through the \texttt{Endomap.power} pathway --- exercising the
binary powering loop, not a naive iteration --- and asserts idempotence of
the result by extensional comparison, the executable face of the safe
finite power theorem of Chapter~5. The family includes the degenerate
sizes ($0! = 1! = 1$), all permutations (where $f^{|X|!} = \id$), all
constants, and every mixed tail-and-cycle shape through four points.

The choice of bound four is not arbitrary, and it is worth recording what
size four buys that size three does not. The proof of the factorial
theorem (Chapter~5) has two moving parts --- tail absorption and cycle
divisibility --- and the interesting stress cases are those where the
least idempotent exponent is governed by the \emph{least common multiple}
of several cycle lengths rather than by any single cycle. On three points
the possible cycle structures are limited (cycle lengths $1$, $2$, $3$,
never two nontrivial cycles at once); on four points the double
transposition appears --- two disjoint $2$-cycles --- along with the
$4$-cycle, the $3$-cycle-plus-fixed-point, and every combination of a
nontrivial cycle with a genuine tail. Size four is thus the first size at
which the family contains maps whose cycle structure is plural, and
$4! = 24$ is the first factorial that must simultaneously absorb a
length-$3$ tail bound and divide out cycle lengths $1$ through $4$. The
family also stresses the powering \emph{implementation} distinctly at
each size: the exponents actually exercised are $1, 1, 2, 6, 24$, whose
binary expansions ($1$, $10$, $110$, $11000$) drive different
accumulate/square interleavings through the loop of Chapter~9.

\emph{Typed maps through size four: $499$ maps.} Across ordered pairs of
carrier sizes $(m, n)$ with $0 \le m, n \le 4$, the total maps from an
$m$-element source to an $n$-element target number $n^m$ --- zero when
$n = 0 < m$, one when $m = 0$. Summing:

\[
\sum_{m=0}^{4}\sum_{n=0}^{4} n^{m}
\;=\; 5 + 10 + 30 + 100 + 354 \;=\; 499 ,
\]

again a derivation consistent with the recorded count (the five terms are
the sums for $m = 0, 1, 2, 3, 4$ respectively). These typed families
drive the constructor validation and the composition guard across
\emph{distinct} carriers: identity neutrality with typed identities on
both sides, and rejection behavior --- every source/target mismatch
raising \texttt{CompositionTypeError} --- are checked here, where
same-carrier coincidences cannot mask a name-blind or cardinality-blind
comparison.

The suite's regression (as opposed to exhaustive) component pins the
powering machinery at scale: the exponents $2^{30}$ and $10^{5000}$
exercise the bounded presentation-name construction of Chapter~9 ---
the second exponent would overflow Python's decimal-conversion limit if
names were built naively --- and the negative fixtures drive sets,
mappings, generators, blank names, out-of-range indices, boolean indices,
and negative exponents into their respective \texttt{ValidationError} and
\texttt{CompositionTypeError} rejections, one assertion per taxon of the
error taxonomy.

\subsection{Lift structure and functor-law regressions}\label{ch10:sec:lift-suite}

The lift suite is the largest test module by byte count, and its central
family is the exact functoriality sweep.

\emph{Lift-functor composition pairs through size three: $1{,}678$
pairs.} Consider all composable typed pairs $f : X \to Y$, $g : Y \to Z$
with $|X|, |Y|, |Z| \le 3$, including empty carriers. For fixed sizes
$(a, b, c)$ there are $b^a \cdot c^b$ such pairs, and summing over the
middle size first,

\[
\sum_{b=0}^{3}\Bigl(\sum_{a=0}^{3} b^{a}\Bigr)\Bigl(\sum_{c=0}^{3} c^{b}\Bigr)
\;=\; 1\cdot 4 + 4\cdot 6 + 15\cdot 14 + 40\cdot 36
\;=\; 4 + 24 + 210 + 1{,}440 \;=\; 1{,}678 ,
\]

a derivation consistent with the manifest count. For every pair the suite
builds the three matrices $A_g$, $A_f$, and $A_{g \circ f}$ through
\texttt{DeterministicLift.matrix()}, multiplies $A_g A_f$ in exact integer
arithmetic, and asserts \emph{entrywise equality} --- tuples of tuples of
\texttt{int}s compared with \texttt{==}, no tolerance, no floating point
anywhere in the comparison. Identity preservation $A_{\id_X} = I$ is
checked for every carrier in scope by the same exact comparison. This is
the executable half of the lift functor laws (Chapter~7); the
kernel-checked half, and the gap between them, is Chapter~11's subject.

The structural regressions of the module tie the theorem-as-shortcut
properties of \texttt{deterministic\_lift.py} back to matrix reality on
the bounded family: fiber sizes against explicitly counted preimages, the
squared norm against the maximum fiber size, the trace shortcut against
the actual Gram diagonal sum, rank against the count of nonempty fibers,
the input Gram matrix against a directly computed $A_f^{\mathsf T} A_f$,
the output Gram against the diagonal of fiber sizes, and the support
projector against the zero-one support pattern --- with fixtures where
$A_f A_f^{*}$ and the support projector \emph{differ} (any fiber of size
two or more), so the three-operator distinction of Chapter~9 is enforced
by test, not just by method naming. The permanent noncontractivity fixture
--- the constant map on two points with $A_f^2 = A_f$, norm $\sqrt 2$,
Gram trace $2$, rank $1$, neither a partial isometry nor an orthogonal
projector --- is asserted property by property, standing guard over the
two tempting over-readings (``deterministic lifts are contractive,''
``idempotent lifts are projectors'') that Chapter~7 identified and the
theorem ledger records as \textsc{thm-lift-occupation-counterexample}.

\subsection{Normal-form correctness and resource regressions}\label{ch10:sec:nf-suite}

The normal-form suite verifies the closure engine of Chapter~9 on
hand-checkable catalogs, in three bands. \emph{Correctness:} for small
generator catalogs whose closures can be enumerated by hand (of the kind
worked in full in Appendix~A), the suite asserts the exact denotation
count, the exact representative word for each reachable denotation ---
shortlex-least, verified against independent enumeration --- the identity
of \texttt{normalize} on words with equal denotations, and the stability
of the \texttt{normal\_forms} listing order. \emph{Rejection:}
constructor-level negative fixtures cover duplicate generator labels,
duplicate object names, generators referencing undeclared objects, unknown
labels at evaluation, undeclared start objects, and unordered container
inputs; \texttt{representative\_for} on an arrow outside the generated
closure is asserted to raise \texttt{UnreachableDenotationError}
specifically. \emph{Resource guards:} catalogs constructed with
deliberately small \texttt{max\_arrows} or \texttt{max\_frontier} are
asserted to raise \texttt{ClosureLimitExceeded} --- and, decisively, to
leave \emph{no} accessible partial catalog: the exception propagates from
the constructor, so there is no object on which a truncated table could
be observed. The fail-closed policy is thereby a tested behavior, not a
documentation promise.

\subsection{The manifest-boundary regression}\label{ch10:sec:boundary-suite}

The fourth module is small --- under two kilobytes --- and unusual in
kind: it tests the \emph{claim discipline} rather than the code. It loads
the prototype manifest and asserts the integration boundary fields on
which every downstream reading of the artifact depends:

\begin{itemize}
\item \texttt{active\_cas\_compiler\_modified} is \texttt{false} --- the
      kernel has not touched the production notation compiler;
\item \texttt{production\_runtime\_dependency} is \texttt{false} --- no
      production system depends on the kernel;
\item \texttt{canonical\_tks\_v8\_status} is \textsc{open} --- the kernel
      is not canonical TKS v8;
\item \texttt{cas\_semantic\_transfer\_status} is
      \textsc{open / not licensed} --- no CAS semantic content has been
      mapped onto the kernel.
\end{itemize}

\begin{remark}[Why boundaries belong in the suite]\label{ch10:rem:boundary}
A boundary that lives only in prose can drift silently: a manifest edit
that flipped \texttt{cas\_semantic\_transfer\_status} to a permissive
value would, without this module, leave the suite green while the
artifact's most important commitment quietly inverted. Placing the
assertions in the suite makes boundary drift a \emph{test failure}: the
20-test green result certifies not only that the code behaves but that
the artifact still claims exactly as little as it is licensed to claim.
This is the same instinct as the permanent noncontractivity fixture ---
encode the prohibition where the machinery will trip over it --- applied
to the artifact's governance metadata.
\end{remark}

\section{Coverage summary}\label{ch10:sec:coverage}

\begin{center}
\begin{tabular}{@{}lrl@{}}
\toprule
Bounded exhaustive family & Cases & Law exercised \\
\midrule
Three-point associativity triples & $19{,}683$ & associativity (Ch.~4) \\
Endomaps on carriers through size four & $289$ & factorial idempotence (Ch.~5) \\
Typed maps, sizes through four & $499$ & typing, identity, rejection (Ch.~4) \\
Lift composition pairs through size three & $1{,}678$ & functor laws, exact (Ch.~7) \\
\bottomrule
\end{tabular}
\end{center}

All four counts are recorded in the manifest's validation snapshot
(\texttt{validated\_at} 2026-08-04, \texttt{result} \textsc{pass},
\texttt{tests\_run} 20) and re-derived arithmetically above. The
per-family scopes differ deliberately: associativity burns its budget on
one carrier at depth three (the law is about composition structure, not
carrier variety); the endomap and typed-map families spend theirs on
carrier variety through size four (the laws are about map shapes); the
functor sweep stops at size three because its cost is cubic in the three
sizes jointly and the matrix shapes it must distinguish --- collision,
non-surjectivity, empty domain and codomain --- all occur by then.

Two structural remarks on the families' composition. First, the empty
cases are inside the bounds, not excluded from them: the typed-map family
includes the $m = 0$ column (one empty map into each target, including
the empty map on the empty object) and the $n = 0 < m$ cells (no maps at
all, which the enumerator must produce as an empty family rather than an
error), and the lift sweep explicitly includes empty carriers ---
zero-by-$n$ and $n$-by-zero matrices whose products are degenerate in
ways that off-by-one code notoriously mishandles. The empty cases carry
disproportionate defect-finding weight for exactly that reason. Second,
the four families overlap by construction --- a three-point endomap
appears in the associativity family, the endomap family, the typed-map
family, and (lifted) in the functor sweep --- and the overlap is a
feature: the same mathematical object is exercised through four
different code paths (composition, powering, construction and typing,
matrix building), so an inconsistency \emph{between} code paths, not
just within one, has four chances to surface.

A note on cost. The largest family runs $19{,}683$ triple compositions
of three-element mappings; the functor sweep performs $1{,}678$ exact
$3 \times 3$-bounded integer matrix multiplications; the whole
20-test suite completes in seconds on commodity hardware. The
bounded-exhaustive design's practical thesis is visible in these
numbers: at kernel scale, exhaustion within the interesting scope is
cheaper than the infrastructure of sampling would be. The suite spends
its complexity budget on \emph{which} families to exhaust --- the
design decisions recorded above --- and none on test-data machinery.

\section{Trusting the tests}\label{ch10:sec:trust}

A verification methodology that scrutinizes the code and exempts the
tests has merely relocated its trust, not reduced it. Three mechanisms
address the suite's own trustworthiness; none eliminates it, and the
residue is stated at the end.

\emph{First, the tests are bound.} All four test modules carry SHA-256
bindings in the manifest with declared roles, and the Gate~1 verifier
hashes them before running them (Section~\ref{ch10:sec:gate1}). A
weakened assertion, a deleted family, or a broadened exception match is a
digest change, detectable against the manifest exactly as an
implementation edit would be. Test code is artifact, not scaffolding.

\emph{Second, the coverage counts are externally derivable.} Each
exhaustive family's cardinality is a closed-form combinatorial quantity
--- $(3^3)^3$, $\sum_k k^k$, $\sum_{m,n} n^m$,
$\sum_b (\sum_a b^a)(\sum_c c^b)$ --- recomputed independently in this
chapter and checkable by any reviewer with a pencil. A suite that
silently enumerated less than it claimed would disagree with arithmetic,
not merely with its own comments. This is a modest but real defense: it
pins the \emph{quantifier} of each exhaustive claim to something outside
the code.

\emph{Third, the assertion style is adversarial to weakening.} The
negative fixtures assert specific exception types
(\texttt{CompositionTypeError}, not bare \texttt{Exception}) so a
refactor that degraded a typed rejection into a generic crash would
fail the suite; the fixture tests assert exact values (norm squared
\emph{equals} $2$), not inequalities; and the boundary regression
asserts exact manifest field values, not field presence.

The residue is honestly statable: the suite cannot certify its own
\emph{sufficiency}. Nothing in the mechanisms above excludes the
possibility that some behavior of the implementation is simply never
exercised --- a code path outside every family and every fixture ---
and no coverage measurement is recorded in the manifest snapshot.
Sufficiency review is part of what independent specialist review
(PO-TKS8-CORE-10) is for, and the suite's self-discipline is a floor
under that review, not a substitute for it.

\section{The Gate~1 verifier}\label{ch10:sec:gate1}

\subsection{Protocol}

The script \texttt{scripts/verify\_tks\_v8\_gate1.py} --- itself
hash-bound in the manifest (4{,}230 bytes, role ``hash and executable test
verifier'') --- composes the two evidence mechanisms into a single
replayable check, invoked as
\texttt{python scripts/verify\_tks\_v8\_gate1.py} from the repository
root.

\begin{protocol}[Gate~1 verification]\label{ch10:proto:gate1}
The verifier performs, in order:
\begin{enumerate}
\item \emph{Hash check.} For every artifact binding in the prototype
      manifest, recompute the SHA-256 digest of the file at the bound path
      and compare it, together with the byte count, against the recorded
      values. Any missing file or digest mismatch fails the run before
      any test executes.
\item \emph{Suite rerun.} Execute the four-module test suite and require
      the full pass.
\end{enumerate}
A successful run therefore certifies a conjunction: \emph{these exact
bytes} --- and not some edited neighbor of them --- \emph{pass these
exact tests}.
\end{protocol}

The ordering matters. Running hashes first means a green verifier run can
never be produced by testing modified code: the failure mode ``the tests
pass because someone weakened them'' is caught at step one, since the test
modules are themselves bound. Conversely, the failure mode ``the hashes
match but the code is broken'' is caught at step two. The two steps are
individually cheap and jointly close the obvious substitution attacks on
either evidence type alone --- with the single documented exception, the
manifest's own unbound digest, discussed in Chapter~9: until a reviewed
commit binds the manifest hash, a consistent rewrite of manifest plus
files together would pass. The verifier inherits that boundary; it does
not resolve it.

\subsection{Relation to the Lean verifier}

A second verifier, \texttt{scripts/verify\_psr8\_lean\_gate1.py},
reproduces the formal seed under the pinned Lean toolchain; it is the
subject of Chapter~11 and is deliberately separate: executable evidence
and kernel-checked evidence have different trust bases, different failure
modes, and different scopes, and conflating them in one green light would
blur exactly the distinctions this monograph exists to keep sharp. Gate~1
reproduction means running both.

\section{What a green suite does not establish}\label{ch10:sec:negative}

This section is normative for every consumer of the artifact, and its
content is quoted or directly derived from the manifest's own boundary
field: ``A passing bounded executable suite is not proof-assistant
verification or independent specialist review.''

\begin{boundary}[Non-claims of the executable evidence]\label{ch10:bdy:negative}
A full pass of the 20-test suite, including under the Gate~1 verifier,
does \emph{not} establish any of the following.
\begin{enumerate}
\item \emph{Universal correctness.} The exhaustive families are bounded at
      carrier sizes three and four. The suite proves nothing about size
      five --- the universal statements are carried by the proofs of
      Chapters~4--8, and their formalization status is the open ledger of
      Chapter~12, not a consequence of testing.
\item \emph{Proof-assistant verification of the Python.} The Lean seed of
      Chapter~11 checks a Lean model of the definitions, not the Python
      source; equivalence between the two is itself an open obligation.
      A green suite plus a green Lean build still leaves the running code
      unverified in the proof-assistant sense.
\item \emph{Independent review.} All tests were authored within the
      project. PO-TKS8-CORE-10 (independent specialist review) has no
      receipt and remains \textsc{open}; nothing a self-authored suite
      does can close it.
\item \emph{Closure of any open obligation.} The ledger rows marked
      \textsc{open} in Chapter~12 --- formalizations, compiler
      integration, semantic reconciliation, review --- are untouched by
      any number of green runs; the suite's role in the ledger is
      confined to the rows explicitly marked \textsc{executable
      validation}.
\item \emph{Semantic adequacy for TKS/CAS.} The suite contains no Noetic,
      Domain, Drive, or Register content --- necessarily, since the code
      contains none. Green tests say nothing about whether the kernel's
      types are the \emph{right} substrate for the canonical ontology;
      that is PO-TKS8-CORE-08, \textsc{open / not licensed}.
\item \emph{Empirical claims of any kind.} No cognitive-dynamics,
      channel-contractivity, or CAS-efficacy claim follows from any test
      in the suite, per the standing claim boundaries of the manifest.
\end{enumerate}
\end{boundary}

\begin{remark}[The suite as a floor, not a ceiling]\label{ch10:rem:floor}
The correct reading of the executable evidence is as a floor under the
artifact: the definitions are implementable, the implementation agrees
with the proved mathematics everywhere a bounded check can look, the
rejection behavior matches the error taxonomy case by case, the
fail-closed guards actually fail closed, and the claim boundaries are
mechanically pinned. Everything above that floor --- universality,
formal verification of the running code, review, semantic transfer ---
is exactly as open as the ledger says it is. Keeping the floor and the
ceiling distinct is the methodology.
\end{remark}
